There is a positive, quantified sub-contract. Change the native row
\[
 66/0:\quad 1\,R\,198\quad\longmapsto\quad 0\,L\,54
\]
and delete native state 198 (dense state 189). The resulting research table has
214 listed working labels and 428 rows; 212 working labels remain reachable
in its transition graph. The inherited symbolic runtime checker proves
all 1,900 literal increment occurrences with the exact logical increment and
continuation, while leaving the arbitrary nonempty spare suffix unchanged.
Its induction rules cover arbitrary natural register values, arbitrary natural
gap and arbitrary finite spare data; this is stronger than bounded testing.
The run uses 672,539 symbolic rules. It is an ordinary quantified symbolic
argument checked by a Python verifier, not a proof-assistant kernel theorem.

The same table fails a concrete zero-decrement contract. At zero-based source
index 2 (the third leaf),
address \texttt{0100}, command \texttt{CD0}, take all nine registers zero,
spare suffix $[0]$ and gap zero. The parent returns after 640 transitions with
all logical registers zero and continuation one-bits at positions $[1,2]$.
The research table first returns after 3,966 transitions with continuation
one-bits at $[1,3,4,5]$ and register values
\[
 [0,0,0,0,0,1,0,0,2].
\]
An independent simulator replays this mismatch. Thus the 214-state table is
delivered only as a rejected uniform-runtime experiment, with its failed
obligation prominently identified. It is not delivered as a verified solution
to the consistency machine problem; its blank-run calibration remains
unestablished.

All 856 direct write/direction/target replacements for the incoming row also
fail the stated relaxed short screen. A targeted 856-choice provider experiment
likewise finds no sampled full-contract survivor. These are bounded
observations, not global impossibility results for a joint rewrite of the
shared region.

\subsection{Coordinated reuse of an existing provider}

A further fixed family redirects dense row $66/0$ to another provider $p$ and
copies either or both outgoing rows of state 189 into $p$, then deletes 189.
Excluding providers 66 and 189 gives $213\cdot3=639$ tables. All 109 survivors
of the short screen fail the subsequent complete-source sample refinement.
An independent simulator audits the whole family, with the following results.

\begin{center}
\begin{tabular}{lr}
\toprule
Observed obstruction or remaining bound & Tables\\\midrule
Wrong first-redispatch output & 338\\
Outward motion into the infinite blank exterior & 223\\
Literal repeated configuration & 5\\
Unexpected halt during a nonhalting primitive & 2\\
Unresolved at the 20,000-transition bound & 71\\\bottomrule
\end{tabular}
\end{center}

Thus 568 tables have explicit primitive-contract obstructions; the remaining
71 have only bounded failures. The ray checks verify a read-zero self-loop
moving away from every nonzero cell, and the cycle checks compare complete
configurations. These conclusions are stronger than timeouts but remain
specific to the stated primitive contract.

Three variants, using providers 95, 188 and 213 with both rows replaced, are
particularly informative. Each has 214 listed and graph-reachable working
states. Their first decoded blank-start frames match the parent's complete
tape and head: at 79,633 transitions for the first two, and 79,639 for the
third. This experimental first decoded frame is distinct from the maintained
formal bootstrap endpoint. They pass respectively 5,944, 5,936 and 5,918
guarded source cases before their first failed case. None of those finite
successes establishes correctness.

Their failures occur at zero-based source indices 4,578, 4,572 and 4,560.
The bad prefixes read only positions $-10$ through 14, 15 or 21 relative to
the initial head. Every input cell depending on registers, gap or spare data
lies at position at least 64. The returned continuation differs from both
branches of the source decrement. The same finite prefix therefore fails for
every natural register file and every legal gap and spare suffix at that
address. A rescue by a stronger input invariant must exclude the address
itself; excluding only the displayed zero-register valuation is insufficient.
The independently replayed tables and witnesses are retained as rejected
experiments in \texttt{state\_search/forced\_provider189/}.

The addresses also have a useful source-level explanation. They lie in the
mandatory final gate of the AX routine: the first comparison, restoration of
its saved value, and the second comparison. Direct graph replay reaches a
\texttt{Clean} AX entry after 180 checkpoint transitions from the supplied
initial decoded frame. The inherited total macro-run theorems then reach the
gate with registers $[7,T,W,0,0,0,0,0,0]$ for finite $T,W$. The literal graph
visits source indices 4,560, 4,572 and 4,578 at offsets $1$, $T+4$ and $T+6$
from gate entry. Thus an ordinary argument from inherited totality and the
checked literal mapping establishes their eventual occurrence in the parent
source run; a long bounded trace was merely still doing earlier arithmetic.
This argument is documented in
\texttt{state\_search/source\_reachability/}. It is not a new kernel
reachability theorem, and it does not establish that any altered physical
machine follows the parent to those controls or has a different blank-tape
halting behavior.

The overwritten providers are source-address decoder controls. Provider 188
is especially localized: it separates two closed transfer loops, from
register 7 to 6 and from register 4 to 5. Three simple swaps or reassociations
of this region preserve 202 distinct source DAG nodes; the corresponding
literal two-node rewires each still have 215 working states and no ordinary
Mealy-state merge. The source transformations are justified by closed-block
sequencing and the disjoint register pairs. They do not certify those physical
rewires: a changed intermediate decoder context still needs its own runtime
proof. This small experiment saves no state. A more substantial joint
implementation of the two transfers, their decoder and the shared
append/restoration interface remains a concrete synthesis target.
